\documentclass{beamer}

\input{../helper}

\title{Introdução à Computação \\ escopo de variáveis e  arrays}
\author{Alexandre Rademaker}
\date{}

\begin{document}

\begin{frame}
 \maketitle
\end{frame}


\begin{frame}[fragile,allowframebreaks]{variáveis e escopo}

  Escopo é uma característica de uma variável que define a partir de
  quais funções essa variável pode ser acessada.
  
  As variáveis \emph{locais} só podem ser acessadas dentro das funções nas
  quais são criadas.

  Variáveis \emph{globais} podem ser acessadas por qualquer função no
  programa.

  \framebreak

  \begin{lstlisting}[style=normalc]
    int triple(int x)
    {
      return x * 3;
    }
    
    int main(void)
    {
      int result = triple(5);
    }

  \end{lstlisting}

  \ilcode{x} é local na função \ilcode{triple} e \ilcode{result} é
  local na \ilcode{main}.

  \framebreak

  \begin{lstlisting}[style=normalc]
    #include <stdio.h>
    
    float global = 0.5050;

    void triple(void)
    {
      global *= 3;
    }
    
    int main(void)
    {
      triple();
      printf("%f\n", global);
    }
  \end{lstlisting}

  \framebreak

  Na maioria das vezes, as variáveis locais em C são passadas por
  valor nas chamadas de função.
  
  Quando uma variável é passada por valor, a função \emph{chamada}
  recebe uma cópia da variável passada, não a própria variável.

  Isso significa que a variável no \emph{chamador} permanece
  inalterada, a menos que seja substituída.

  \framebreak
  
  O que será impresso? 4 ou 12?
  
  \begin{lstlisting}[style=ssc]
    int triple(int x) {
      return x *= 3;
    }
    
    int main(void) {
      int foo = 4;
      triple(foo);
      printf("%d", foo);
    }
  \end{lstlisting}

  \begin{lstlisting}[style=ssc]
    int triple(int x) {
      return x *= 3;
    }
    
    int main(void) {
      int foo = 4;
      foo = triple(foo);
      printf("%d", foo);
    }
  \end{lstlisting}

  \framebreak

  As coisas podem ficar particularmente enganadora se o mesmo nome de
  variável aparecer em várias funções, o que é perfeitamente
  aceitável, desde que as variáveis existam em escopos diferentes.

  \framebreak

  \begin{lstlisting}[style=normalc,escapechar={|}]
    int main(void)
    {
      int |\color{cyan}x |= 1;
      int y;
      y = increment(|\color{cyan}x|);
      printf("x is %i, y is %i\n", |\color{cyan}x|, y);
    }
    
    int increment(int |\color{orange}x|)
    {
      |\color{orange}x|++;
      return |\color{orange}x|;
    }
  \end{lstlisting}

  {\centering x is 1, y is 2}
  
\end{frame}


\begin{frame}[fragile,allowframebreaks]{arrays}

  Um array é um bloco de espaço contíguo na memória
  
  que foi particionado em pequenos blocos de espaço de tamanho
  idêntico chamados elementos
  
  cada um dos quais pode armazenar uma certa quantidade de dados
  
  todos do mesmo tipo de dados, como \ilcode{int} ou \ilcode{char}

  e que pode ser acessado diretamente por um índice.

  \framebreak

  Em C, os $n$ elementos de um array são indexados a partir de 0. O
  último elemento estará na posição $n-1$

  C é muito indulgente. Nada irá prevenir que você \emph{saia dos
    limites} de seu array; tome cuidado!

  \framebreak

  Declaração

  {\texttt {\color{green} type} {\color{red} name}[{\color{blue} size}];}

  Exemplo

  {\texttt {\color{green} int} {\color{red} student\_grades}[{\color{blue} 30}];}

  \framebreak

  você pode realizar as mesmas operações com um elemento do array que
  faria com variáveis do mesmo tipo.\vspace{.5cm}

  \begin{lstlisting}[style=normalc]
    bool truthtable[10];
    
    truthtable[2] = false;
    
    if(truthtable[7] == true) {
      printf("TRUE!\n");
    }
    truthtable[10] = true;
  \end{lstlisting}

  \framebreak

  Para declarar e inicializar um array, temos um ``açucar sintático''!\vspace{.5cm}

  \begin{lstlisting}[style=ssc,escapechar={|}]
    // instantiation syntax
    bool truthtable|\color{red}[]| = { false, true, true };

    // individual element syntax
    bool truthtable[3];
    truthtable[0] = false;
    truthtable[1] = true;
    truthtable[2] = true;
  \end{lstlisting}

  \framebreak

  Arrays podem ter mais do que uma dimensão!

  \begin{lstlisting}[language=C,escapechar={|},basicstyle=\ttfamily]
    bool battleship[10][10];
  \end{lstlisting}

  Na memória, de fato, apenas um array de 100 dimensões.

  \framebreak

  Embora possamos tratar elementos individuais dos arrays como
  variáveis, não podemos tratar arrays inteiros como variáveis.
  
  Não podemos, por exemplo, atribuir um array a outro usando o
  operador de atribuição.

  Em vez disso, devemos usar um loop para copiar os elementos um de
  cada vez.

  \begin{lstlisting}[language=C,escapechar={|},basicstyle=\ttfamily]
    int foo[5] = { 1, 2, 3, 4, 5 };
    int bar[5];
    |{\color{red} \st{bar = foo;}|
  \end{lstlisting}

  \framebreak

  \begin{lstlisting}[language=C,escapechar={|},basicstyle=\ttfamily]
    int foo[5] = { 1, 2, 3, 4, 5 };
    
    int bar[5];
    for(int j = 0; j < 5; j++)
    {
      bar[j] = foo[j];
    }    
  \end{lstlisting}

  \framebreak

  Lembre-se que variáveis em C são passadas \emph{por valor} nas chamadas de função.
  
  Arrays não seguem esta regra. Eles são passados \emph{por
    referência}. A função chamada recebe o array, não uma cópia dele.
  
  O que isso significa quando a função chamada manipula elementos do array?
  
  Por enquanto, apenas vamos entender que arrays têm essa propriedade
  especial, mas voltaremos ao assunto em breve!

  \framebreak

  \begin{lstlisting}[style=ssc]
    void set_array(int array[4]) {
      array[0] = 22;
    }
    
    void set_int(int x) {
      x = 22;
    }
    
    int main(void)
    {
      int a = 10; int b[4] = { 0, 1, 2, 3 };
      set_int(a);
      set_array(b);
      printf("%d %d\n", a, b[0]); // prints 10, 22
    }
  \end{lstlisting}

\end{frame}


\end{document}

%%% Local Variables:
%%% mode: latex
%%% TeX-master: t
%%% End:
